\documentclass[12pt,a4paper]{article}
\usepackage[utf8]{inputenc}
\usepackage[spanish]{babel}
\usepackage{amsmath}
\usepackage{amsfonts}
\usepackage{amssymb}
\usepackage{geometry}
\geometry{a4paper, margin=2.5cm}

\title{Álgebra Lineal: Tema 3 - Espacios Vectoriales}
\author{Explicación detallada}
\date{}

\begin{document}

\maketitle

\section{Introducción: ¿Qué es un Espacio Vectorial?}
Un espacio vectorial es un conjunto no vacío $V$ cuyos elementos se denominan \textbf{vectores}, definido sobre un cuerpo $K$ (como los reales $\mathbb{R}$ o complejos $\mathbb{C}$) cuyos elementos se denominan \textbf{escalares}.

Para ser un espacio vectorial, el conjunto debe permitir dos operaciones (suma interna y producto por un escalar) y verificar \textbf{8 propiedades}:

\textbf{Propiedades de la suma (Operación interna):}
\begin{itemize}
    \item \textbf{Asociativa:} $(u+v)+w = u+(v+w), \forall u,v,w \in V$.
    \item \textbf{Conmutativa:} $u+v = v+u, \forall u,v \in V$.
    \item \textbf{Elemento neutro:} $\exists 0 \in V \mid 0+v = v, \forall v \in V$.
    \item \textbf{Elemento opuesto:} $\forall v \in V, \exists -v \in V \mid v + (-v) = 0$.
\end{itemize}

\textbf{Propiedades del producto por un escalar (Operación externa):}
\begin{itemize}
    \item \textbf{Distributiva I:} $\alpha \cdot (u+v) = \alpha \cdot u + \alpha \cdot v$.
    \item \textbf{Distributiva II:} $(\alpha+\beta) \cdot v = \alpha \cdot v + \beta \cdot v$.
    \item \textbf{Asociativa:} $\alpha \cdot (\beta \cdot v) = (\alpha\beta) \cdot v$.
    \item \textbf{Elemento unidad:} $\exists 1 \in K \mid 1 \cdot v = v, \forall v \in V$.
\end{itemize}

\textbf{Ejemplos clásicos}: 
$\mathbb{R}^n$, las matrices $\mathcal{M}^{m\times n}(\mathbb{R})$, los polinomios $\mathcal{P}_n(\mathbb{C})$, y el espacio trivial $\{0\}$[cite: 51, 54, 57, 61].

\section{Dependencia e Independencia Lineal}
\begin{itemize}
    \item \textbf{Combinación Lineal (c.l.):} Es un vector generado mediante $\alpha_1 v_1 + \dots + \alpha_n v_n$.
    \item \textbf{Dependencia Lineal (l.d.):} Un vector es l.d. si es una c.l. del resto. Un conjunto es l.d. si existen escalares no todos nulos tales que $\alpha_1 v_1 + \dots + \alpha_n v_n = 0$. Si un conjunto es l.d., cualquier superconjunto también lo es. Es l.d. si al menos un vector se puede expresar como c.l. del resto.
    \item \textbf{Independencia Lineal (l.i.):} Si $\alpha_1 v_1 + \dots + \alpha_n v_n = 0$ implica obligatoriamente que $\alpha_1 = \dots = \alpha_n = 0$. Todo subconjunto de un conjunto l.i. sigue siendo l.i.
\end{itemize}

\textbf{Cálculo mediante el Rango:}
Para los vectores columnas en la matriz $(v_1|\dots|v_n)$:
\begin{itemize}
    \item $l.i. \Leftrightarrow \text{S.C.D.} \Rightarrow rg(v_1|\dots|v_n) = n$.
    \item $l.d. \Leftrightarrow \text{S.C.I.} \Rightarrow rg(v_1|\dots|v_n) < n$.
\end{itemize}
El rango del conjunto es el número de vectores que son l.i.

\section{Base y Dimensión}
\begin{itemize}
    \item \textbf{Sistema Generador (s.g.):} Conjunto $S$ tal que cualquier vector del espacio se puede expresar como c.l. de los elementos de $S$.
    \item \textbf{Base ($\mathcal{B}$):} Es un conjunto de vectores que es linealmente independiente y además es sistema generador.
    \item \textbf{Dimensión ($dim(V)$):} Es el número de vectores que conforman una base. Por el Teorema principal, todas las bases de un mismo espacio vectorial $V$ tienen el mismo número de elementos[cite: 150, 151].
\end{itemize}

Si sabemos que $dim(V)=n$ y $S \subset V$:
\begin{itemize}
    \item Si $S$ es l.i. $\Rightarrow card(S) \leq n$.
    \item Si $S$ es s.g. $\Rightarrow card(S) \geq n$.
    \item Si $S$ es l.i. y $card(S)=n \Rightarrow$ es Base.
    \item Si $S$ es s.g. y $card(S)=n \Rightarrow$ es Base.
\end{itemize}

\section{Coordenadas respecto de una Base}
Sea $\mathcal{B}=\{v_1, \dots, v_n\}$ una base. Todo vector $v \in V$ se puede expresar de \textbf{forma única} como c.l. de $\mathcal{B}$:
$$v = \alpha_1 v_1 + \dots + \alpha_n v_n$$ [cite: 198, 200]
A los coeficientes se les llama \textbf{coordenadas} respecto a $\mathcal{B}$, denotado como $v(\alpha_1, \dots, \alpha_n)_{\mathcal{B}}$. Si el subíndice se omite, se asume la base canónica. Las coordenadas de los vectores de la base respecto a ella misma conforman la matriz identidad.
Un conjunto de vectores es l.i. si y solo si el rango de la matriz formada por sus coordenadas en la base $\mathcal{B}$ es igual a $n$.

\section{Cambio de base}
Dadas dos bases $\mathcal{B}$ y $\mathcal{B}'$. 
Para cambiar las coordenadas de un vector de la antigua base $\mathcal{B}$ a la nueva $\mathcal{B}'$, usamos la ecuación:
$$x_{\mathcal{B}'} = P_{\mathcal{B}}^{\mathcal{B}'} x_{\mathcal{B}}$$ [cite: 256, 257]
Las columnas de la matriz de paso $P$ son las coordenadas de los vectores de la base antigua $\mathcal{B}$ expresados en la base nueva $\mathcal{B}'$[cite: 236, 237, 259].

Para el viaje inverso (de $\mathcal{B}'$ a $\mathcal{B}$), se utiliza la matriz inversa:
$$x_{\mathcal{B}} = (P_{\mathcal{B}}^{\mathcal{B}'})^{-1} x_{\mathcal{B}'}$$ 
Donde las columnas de la matriz inversa representan las coordenadas de la base nueva $\mathcal{B}'$ en la base antigua $\mathcal{B}$.

\end{document}